\chapter{Introduction}
\label{ch:intro}

% ===========================================================================
% FRAMING, SETTLED WITH THE SUPERVISOR 2026-09-22.
%
% This is a DESIGN-AND-EVALUATION study, not a hypothesis-confirmation study.
% The motivation was to design an agentic simulation architecture, evaluate it
% honestly against an old-style scripted baseline, and report what was learned
% either way. A negative answer is a result here, not a failure.
%
% Write the whole thesis in that voice. Never apologise for the outcome.
% ===========================================================================

\section{Motivation}
\label{sec:intro:motivation}
% Why evacuation coordination is hard: heterogeneous organisations, local
% information, a hard deadline, and a disruption space too large to enumerate.
% Why the old style falls short: a pre-scripted workflow can only cover the
% cases its author thought of.
% Why an agentic architecture is the obvious thing to try: agents that reason
% over a shared database can in principle answer a disruption nobody scripted.
% END ON THE OPEN QUESTION, not on a promise: nobody had measured whether that
% promise survives contact with a scored evacuation.

\section{Problem Statement}
\label{sec:intro:problem}
% The concrete problem: move 100 residents of a riverside nursing home to
% shelters before a forecast storm surge, using resources owned by eight
% different organisations, when the plan must change mid-operation.

\section{Research Questions}
\label{sec:intro:rq}
% ---------------------------------------------------------------------------
% THE FOUR QUESTIONS, AS SETTLED. Each has an answer in this thesis. Write the
% question here and the answer in the chapter named, then again in the
% conclusion. Do not hedge any of the four.
%
% RQ1  Can a database-centric agentic architecture coordinate an evacuation
%      BETTER than an old-style scripted workflow baseline?
%      ANSWER: No, on the pre-registered primary outcome.
%      -> Chapter~\ref{ch:results}. Gate 1 fails 26 of 27 cells. The baseline
%         saved every resident in every novel run; the best agentic
%         configuration saved 90 %.
%      -> And it is worse than that: against a null policy that ignores every
%         disruption, agentic regret is negative in all 15 novel cells and 68
%         of 150 runs came out WORSE THAN DOING NOTHING.
%      -> BOUNDARY, state it in the same breath (Chapter~\ref{ch:critique}):
%         no scenario in this thesis ever REQUIRED adaptation, so this answers
%         "did it win here", not "can it ever win".
%
% RQ2  If not, why not? Where does the deficit come from?
%      ANSWER: A COMMITMENT failure, not a comprehension failure.
%      -> Chapter~\ref{ch:results} and Chapter~\ref{ch:critique}.
%      -> The agents understood the disruption. They were told the right
%         population in 150 of 150 runs, closed 93.8 % of addressed requests,
%         dispatched the same number of missions in failing and succeeding
%         runs, and used the reserve bus in 100 % of runs.
%      -> What they could not do is carry ONE COMMITTED PLAN to delivery.
%         Losses are quantised at whole bus-loads: of 150 runs, 102 delivered
%         everyone, 19 delivered exactly 50, 18 delivered exactly 90.
%      -> Plan revisions band monotonically against delivery: 90 % full
%         delivery at the minimum revision count, 32 % at six or more.
%         r(plan revisions, S) = -0.408, the strongest non-definitional
%         predictor in the data.
%      -> THE IRONY TO NAME OUT LOUD: re-planning on disruption is the
%         architecture's advertised capability, and on these scenarios the
%         correct number of re-plans was ZERO.
%
% RQ3  What would have to change to make an agentic emergency simulation
%      better?
%      ANSWER: constrain commitment, and test on scenarios that discriminate.
%      -> Chapter~\ref{ch:conclusion}, future work.
%      -> The named experiment: a COMMITMENT-CONSTRAINED VARIANT. Cap re-plans,
%         or forbid cancelling a leg already in flight, run against the same
%         scenarios under the same models. One campaign, no new scenario
%         authoring. It is also what would settle the causal direction, which
%         this thesis explicitly does NOT establish.
%      -> Plus: Gate 0 as a precondition on scenario authoring, a larger k, and
%         more than one SIM_TIME_SCALE.
%
% RQ4  What should be considered when designing an AI agentic emergency
%      simulation?
%      ANSWER: Chapter~\ref{ch:design}. This is the generalisable deliverable
%      and the part that stops being about one nursing home in Hamburg.
%
% ---------------------------------------------------------------------------
% The four engineering sub-questions from the original proposal still hold and
% are answered by the architecture chapters. Keep them, subordinate to RQ1-RQ4:
%   1. Data modelling of actors, states, events, dependencies
%      -> Chapter~\ref{ch:data}
%   2. Integration of heterogeneous organisational agents
%      -> Chapter~\ref{ch:agentic}
%   3. Persistent structures for runtime control AND evaluation
%      -> Sections~\ref{sec:data:eventlog} and \ref{sec:data:evalbackbone}
%   4. Comparative performance, agentic vs workflow
%      -> Chapters~\ref{ch:eval} and \ref{ch:results}

\section{Contributions}
\label{sec:intro:contributions}
% FOUR contributions. State them as claims the thesis actually supports.
%
%  1. ARTEFACT. A database-centric agentic architecture for multi-organisation
%     evacuation coordination: 10 agents across 8 organisations, Postgres as the
%     single source of truth and the active message broker, built and running.
%     Plus a doctrine-grounded workflow baseline that reaches audited parity
%     with it on the shared surface (Chapter~\ref{ch:validity}).
%
%  2. EMPIRICAL. The first scored head-to-head of the two on a common
%     evacuation case: 310 runs, 7 scenarios, pre-registered, deterministically
%     scored. The agentic architecture does NOT outperform the scripted
%     baseline, and underperforms a null policy.
%     (This replaces the proposal's expected finding. Say so plainly.)
%
%  3. DIAGNOSTIC. A measured account of WHY, which is the part a future builder
%     can use: the deficit is commitment, not comprehension, and plan churn is
%     its strongest correlate.
%
%  4. METHODOLOGICAL. The shared append-only event log that makes the
%     comparison fair, auditable and replayable; and Gate 0, a treatment-
%     validity check that catches a disruption which does not discriminate
%     a replanner from a non-replanner BEFORE it is scored.
%
% NOTE. The proposal expected contribution 2 to read "agentic degrades more
% gracefully". It does not. Do not leave the old wording anywhere.

\section{Scope and Delimitations}
\label{sec:intro:scope}
% Out of scope: resident-level agents (residents are a passive count), heavy
% workflow engines (Camunda/Temporal/Airflow), modelling the flood itself,
% real-time citizen apps.
% Bounded by design: one case study, one facility, one city.
% Bounded by what was run: one SIM_TIME_SCALE (8.4), k = 10, three models with
% a complete novel arm. H4 was not tested. See Chapter~\ref{ch:critique}.

\section{Thesis Outline}
\label{sec:intro:outline}
% One paragraph over Chapters~\ref{ch:background} to \ref{ch:conclusion}.
% Point out the two chapters a reader might not expect and why they are there:
% Chapter~\ref{ch:validity} (the comparison has to be admissible before it is
% interesting) and Chapter~\ref{ch:critique} (the evaluation is itself a
% finding).
